% =========================================================
% ppt/src/minggu-13.tex
% Slide pertemuan 13
% Bab XIV — Stack, Prosedur, dan Macro
%
% Sub-CPMK 3.5 | CPMK03 | CPL03
% =========================================================

\documentclass[aspectratio=169,11pt]{beamer}
% =========================================================
% ppt/src/preamble-ppt.tex
% Tema, warna, paket, dan makro bersama seluruh deck slide.
%
% Dipakai oleh setiap berkas minggu-NN.tex dengan pola:
%   \documentclass[aspectratio=169,11pt]{beamer}
%   % =========================================================
% ppt/src/preamble-ppt.tex
% Tema, warna, paket, dan makro bersama seluruh deck slide.
%
% Dipakai oleh setiap berkas minggu-NN.tex dengan pola:
%   \documentclass[aspectratio=169,11pt]{beamer}
%   % =========================================================
% ppt/src/preamble-ppt.tex
% Tema, warna, paket, dan makro bersama seluruh deck slide.
%
% Dipakai oleh setiap berkas minggu-NN.tex dengan pola:
%   \documentclass[aspectratio=169,11pt]{beamer}
%   \input{preamble-ppt}
%
% Palet warna disalin dari buku_ajar/config/styles.tex agar
% identitas visual slide sama dengan buku ajar.
% =========================================================

% ---------------------------------------------------------
% 1. Paket
% ---------------------------------------------------------
\usepackage[T1]{fontenc}
\usepackage[utf8]{inputenc}
\usepackage[indonesian]{babel}
\usepackage{lmodern}
\usepackage{tikz}
\usetikzlibrary{arrows.meta, positioning, calc, shapes.geometric}
\usepackage{booktabs}
\usepackage{array}
\usepackage{amsmath}
\usepackage{amsfonts}
\usepackage{pifont}          % \ding untuk tanda checklist
\usepackage{listings}

% ---------------------------------------------------------
% 2. Palet warna OBE (identik dengan buku)
% ---------------------------------------------------------
\definecolor{obeBlue}{HTML}{1F3864}
\definecolor{obeBlueLight}{HTML}{DCE4F5}
\definecolor{obeTeal}{HTML}{0F6E6E}
\definecolor{obeTealLight}{HTML}{D9F0EE}
\definecolor{obeGreen}{HTML}{2E6B34}
\definecolor{obeGreenLight}{HTML}{DEEFDF}
\definecolor{obeAmber}{HTML}{8A5A00}
\definecolor{obeAmberLight}{HTML}{FDF1D6}
\definecolor{obeMaroon}{HTML}{7A2233}
\definecolor{obeMaroonLight}{HTML}{F7E1E5}
\definecolor{obeGray}{HTML}{4A4A4A}
\definecolor{obeGrayLight}{HTML}{F0F0F0}
\definecolor{codeBack}{HTML}{F7F7F4}
\definecolor{codeFrame}{HTML}{C9C9C0}

% ---------------------------------------------------------
% 3. Identitas mata kuliah (dipakai seluruh deck)
% ---------------------------------------------------------
\newcommand{\NamaMK}{Bahasa Pemrograman Tingkat Rendah}
\newcommand{\NamaMKShort}{Bhs.\ Pemrograman Tingkat Rendah}
\newcommand{\KodeMK}{MWCA010}
\newcommand{\Prodi}{Program Studi Teknik Informatika}
\newcommand{\Fakultas}{Fakultas Teknologi Informasi}
\newcommand{\PerguruanTinggi}{Universitas Bale Bandung}
\newcommand{\DosenPengampu}{Cecep Suwanda, S.Si., M.Kom.}
\newcommand{\SemesterTeks}{Semester Ganjil}

% \PertemuanTeks dan \JudulBab ditetapkan di masing-masing
% berkas minggu-NN.tex sebelum \begin{document}.
\providecommand{\PertemuanTeks}{Pertemuan}

% ---------------------------------------------------------
% 4. Tema beamer
% ---------------------------------------------------------
\usetheme{default}
\setbeamertemplate{navigation symbols}{}

\setbeamercolor{normal text}{fg=black!88}
\setbeamercolor{structure}{fg=obeBlue}
\setbeamercolor{alerted text}{fg=obeMaroon}

% --- judul slide: pita biru penuh lebar ---
\setbeamercolor{frametitle}{fg=white,bg=obeBlue}
\setbeamerfont{frametitle}{series=\bfseries,size=\large}
\setbeamertemplate{frametitle}{%
    \nointerlineskip%
    \begin{beamercolorbox}[%
        wd=\paperwidth,ht=2.9ex,dp=1.5ex,
        leftskip=0.4cm,rightskip=0.4cm]{frametitle}%
        \usebeamerfont{frametitle}\insertframetitle%
    \end{beamercolorbox}%
    \nointerlineskip%
    \hbox{\textcolor{obeTeal}{\rule{\paperwidth}{1.1pt}}}%
}

% --- kaki slide: identitas mata kuliah, pertemuan, nomor ---
\setbeamercolor{footline}{fg=white,bg=obeGray}
\setbeamerfont{footline}{size=\scriptsize,family=\sffamily}
\setbeamertemplate{footline}{%
    \leavevmode%
    \hbox{%
        \begin{beamercolorbox}[wd=0.42\paperwidth,ht=2.1ex,dp=1.1ex,
            leftskip=0.4cm]{footline}%
            \usebeamerfont{footline}%
            \KodeMK\ \textbullet\ \NamaMKShort%
        \end{beamercolorbox}%
        \begin{beamercolorbox}[wd=0.32\paperwidth,ht=2.1ex,dp=1.1ex,
            center]{footline}%
            \usebeamerfont{footline}\PertemuanTeks%
        \end{beamercolorbox}%
        \begin{beamercolorbox}[wd=0.26\paperwidth,ht=2.1ex,dp=1.1ex,
            rightskip=0.4cm]{footline}%
            \usebeamerfont{footline}%
            \insertframenumber\,/\,\inserttotalframenumber%
        \end{beamercolorbox}%
    }%
    \vskip0pt%
}

% --- daftar ---
\setbeamercolor{itemize item}{fg=obeTeal}
\setbeamercolor{itemize subitem}{fg=obeTeal}
\setbeamercolor{enumerate item}{fg=obeBlue}
\setbeamertemplate{itemize item}{\raise0.2ex\hbox{\tiny$\blacksquare$}}
\setbeamertemplate{itemize subitem}{\raise0.2ex\hbox{\scriptsize$\bullet$}}

% --- kotak (pengganti tcolorbox pada buku) ---
\setbeamercolor{block title}{bg=obeBlueLight,fg=obeBlue}
\setbeamercolor{block body}{bg=obeBlueLight!45,fg=black!88}
\setbeamercolor{block title alerted}{bg=obeMaroonLight,fg=obeMaroon}
\setbeamercolor{block body alerted}{bg=obeMaroonLight!45,fg=black!88}
\setbeamercolor{block title example}{bg=obeTealLight,fg=obeTeal}
\setbeamercolor{block body example}{bg=obeTealLight!45,fg=black!88}

% ---------------------------------------------------------
% 5. Makro notasi (sama persis dengan buku)
%
% Catatan: \note TIDAK didefinisikan di sini karena beamer
% sudah memakai \note untuk catatan pembicara. Untuk catatan
% di slide dipakai \catatan.
% ---------------------------------------------------------
\newcommand{\code}[1]{\texttt{#1}}
\newcommand{\asm}[1]{\texttt{#1}}
\newcommand{\reg}[1]{\texttt{#1}}
\newcommand{\flag}[1]{\texttt{#1}}
\newcommand{\dir}[1]{\texttt{#1}}
\newcommand{\file}[1]{\texttt{#1}}
\newcommand{\eng}[1]{\textit{#1}}

\newcommand{\hex}[1]{\texttt{#1h}}
\newcommand{\bin}[1]{\texttt{#1b}}
\newcommand{\dec}[1]{\texttt{#1}}

\newcommand{\important}[1]{\textbf{Penting:} #1}
\newcommand{\catatan}[1]{\textbf{Catatan:} #1}
\newcommand{\term}[1]{\textbf{#1}}

\newcommand{\CPL}[1]{\textbf{CPL#1}}
\newcommand{\CPMK}[1]{\textbf{CPMK#1}}
\newcommand{\SubCPMK}[1]{\textbf{Sub-CPMK #1}}

\newcommand{\checkbox}{\ding{113}\hspace{0.35em}}
\newcommand{\checked}{\ding{51}\hspace{0.35em}}
\newcommand{\crossed}{\ding{55}\hspace{0.35em}}

\newcommand{\flowarrow}{\ensuremath{\rightarrow}}

% ---------------------------------------------------------
% 6. Makro tata letak slide
% ---------------------------------------------------------

% Keterangan sumber materi (ketertelusuran ke buku)
\newcommand{\Sumber}[1]{%
    \vfill
    {\scriptsize\color{obeGray}Sumber: #1\par}%
}

% Frame judul deck: \JudulDeck{<judul bab>}{<pertemuan>}
\newcommand{\JudulDeck}[2]{%
\begin{frame}[plain]
    \vspace{0.25cm}
    {\color{obeBlue}\sffamily\bfseries\small \NamaMK\ (\KodeMK)\par}
    \vspace{0.12cm}
    {\color{obeGray}\sffamily\scriptsize
     \Prodi\ \textbullet\ \Fakultas\ \textbullet\ \PerguruanTinggi\par}
    \vfill
    {\color{obeBlue}\sffamily\LARGE\bfseries #1\par}
    \vspace{0.4cm}
    {\color{obeTeal}\sffamily\large #2\par}
    \vspace{0.45cm}
    {\color{obeTeal}\rule{\textwidth}{1pt}}
    \vfill
    {\color{obeGray}\sffamily\small Dosen Pengampu\par}
    \vspace{0.08cm}
    {\sffamily\bfseries \DosenPengampu\par}
    \vspace{0.2cm}
    {\color{obeGray}\sffamily\small \SemesterTeks\par}
    \vspace{0.25cm}
\end{frame}
}

% Frame pemisah bagian
\newcommand{\PemisahBagian}[1]{%
\begin{frame}[plain]
    \vfill
    \begin{center}
        {\color{obeTeal}\sffamily\small\bfseries BAGIAN\par}
        \vspace{0.35cm}
        {\color{obeBlue}\sffamily\Large\bfseries #1\par}
        \vspace{0.35cm}
        {\color{obeTeal}\rule{0.55\textwidth}{0.9pt}}
    \end{center}
    \vfill
\end{frame}
}

% ---------------------------------------------------------
% 7. Gaya penampil kode
% ---------------------------------------------------------
\input{listing-ppt}

%
% Palet warna disalin dari buku_ajar/config/styles.tex agar
% identitas visual slide sama dengan buku ajar.
% =========================================================

% ---------------------------------------------------------
% 1. Paket
% ---------------------------------------------------------
\usepackage[T1]{fontenc}
\usepackage[utf8]{inputenc}
\usepackage[indonesian]{babel}
\usepackage{lmodern}
\usepackage{tikz}
\usetikzlibrary{arrows.meta, positioning, calc, shapes.geometric}
\usepackage{booktabs}
\usepackage{array}
\usepackage{amsmath}
\usepackage{amsfonts}
\usepackage{pifont}          % \ding untuk tanda checklist
\usepackage{listings}

% ---------------------------------------------------------
% 2. Palet warna OBE (identik dengan buku)
% ---------------------------------------------------------
\definecolor{obeBlue}{HTML}{1F3864}
\definecolor{obeBlueLight}{HTML}{DCE4F5}
\definecolor{obeTeal}{HTML}{0F6E6E}
\definecolor{obeTealLight}{HTML}{D9F0EE}
\definecolor{obeGreen}{HTML}{2E6B34}
\definecolor{obeGreenLight}{HTML}{DEEFDF}
\definecolor{obeAmber}{HTML}{8A5A00}
\definecolor{obeAmberLight}{HTML}{FDF1D6}
\definecolor{obeMaroon}{HTML}{7A2233}
\definecolor{obeMaroonLight}{HTML}{F7E1E5}
\definecolor{obeGray}{HTML}{4A4A4A}
\definecolor{obeGrayLight}{HTML}{F0F0F0}
\definecolor{codeBack}{HTML}{F7F7F4}
\definecolor{codeFrame}{HTML}{C9C9C0}

% ---------------------------------------------------------
% 3. Identitas mata kuliah (dipakai seluruh deck)
% ---------------------------------------------------------
\newcommand{\NamaMK}{Bahasa Pemrograman Tingkat Rendah}
\newcommand{\NamaMKShort}{Bhs.\ Pemrograman Tingkat Rendah}
\newcommand{\KodeMK}{MWCA010}
\newcommand{\Prodi}{Program Studi Teknik Informatika}
\newcommand{\Fakultas}{Fakultas Teknologi Informasi}
\newcommand{\PerguruanTinggi}{Universitas Bale Bandung}
\newcommand{\DosenPengampu}{Cecep Suwanda, S.Si., M.Kom.}
\newcommand{\SemesterTeks}{Semester Ganjil}

% \PertemuanTeks dan \JudulBab ditetapkan di masing-masing
% berkas minggu-NN.tex sebelum \begin{document}.
\providecommand{\PertemuanTeks}{Pertemuan}

% ---------------------------------------------------------
% 4. Tema beamer
% ---------------------------------------------------------
\usetheme{default}
\setbeamertemplate{navigation symbols}{}

\setbeamercolor{normal text}{fg=black!88}
\setbeamercolor{structure}{fg=obeBlue}
\setbeamercolor{alerted text}{fg=obeMaroon}

% --- judul slide: pita biru penuh lebar ---
\setbeamercolor{frametitle}{fg=white,bg=obeBlue}
\setbeamerfont{frametitle}{series=\bfseries,size=\large}
\setbeamertemplate{frametitle}{%
    \nointerlineskip%
    \begin{beamercolorbox}[%
        wd=\paperwidth,ht=2.9ex,dp=1.5ex,
        leftskip=0.4cm,rightskip=0.4cm]{frametitle}%
        \usebeamerfont{frametitle}\insertframetitle%
    \end{beamercolorbox}%
    \nointerlineskip%
    \hbox{\textcolor{obeTeal}{\rule{\paperwidth}{1.1pt}}}%
}

% --- kaki slide: identitas mata kuliah, pertemuan, nomor ---
\setbeamercolor{footline}{fg=white,bg=obeGray}
\setbeamerfont{footline}{size=\scriptsize,family=\sffamily}
\setbeamertemplate{footline}{%
    \leavevmode%
    \hbox{%
        \begin{beamercolorbox}[wd=0.42\paperwidth,ht=2.1ex,dp=1.1ex,
            leftskip=0.4cm]{footline}%
            \usebeamerfont{footline}%
            \KodeMK\ \textbullet\ \NamaMKShort%
        \end{beamercolorbox}%
        \begin{beamercolorbox}[wd=0.32\paperwidth,ht=2.1ex,dp=1.1ex,
            center]{footline}%
            \usebeamerfont{footline}\PertemuanTeks%
        \end{beamercolorbox}%
        \begin{beamercolorbox}[wd=0.26\paperwidth,ht=2.1ex,dp=1.1ex,
            rightskip=0.4cm]{footline}%
            \usebeamerfont{footline}%
            \insertframenumber\,/\,\inserttotalframenumber%
        \end{beamercolorbox}%
    }%
    \vskip0pt%
}

% --- daftar ---
\setbeamercolor{itemize item}{fg=obeTeal}
\setbeamercolor{itemize subitem}{fg=obeTeal}
\setbeamercolor{enumerate item}{fg=obeBlue}
\setbeamertemplate{itemize item}{\raise0.2ex\hbox{\tiny$\blacksquare$}}
\setbeamertemplate{itemize subitem}{\raise0.2ex\hbox{\scriptsize$\bullet$}}

% --- kotak (pengganti tcolorbox pada buku) ---
\setbeamercolor{block title}{bg=obeBlueLight,fg=obeBlue}
\setbeamercolor{block body}{bg=obeBlueLight!45,fg=black!88}
\setbeamercolor{block title alerted}{bg=obeMaroonLight,fg=obeMaroon}
\setbeamercolor{block body alerted}{bg=obeMaroonLight!45,fg=black!88}
\setbeamercolor{block title example}{bg=obeTealLight,fg=obeTeal}
\setbeamercolor{block body example}{bg=obeTealLight!45,fg=black!88}

% ---------------------------------------------------------
% 5. Makro notasi (sama persis dengan buku)
%
% Catatan: \note TIDAK didefinisikan di sini karena beamer
% sudah memakai \note untuk catatan pembicara. Untuk catatan
% di slide dipakai \catatan.
% ---------------------------------------------------------
\newcommand{\code}[1]{\texttt{#1}}
\newcommand{\asm}[1]{\texttt{#1}}
\newcommand{\reg}[1]{\texttt{#1}}
\newcommand{\flag}[1]{\texttt{#1}}
\newcommand{\dir}[1]{\texttt{#1}}
\newcommand{\file}[1]{\texttt{#1}}
\newcommand{\eng}[1]{\textit{#1}}

\newcommand{\hex}[1]{\texttt{#1h}}
\newcommand{\bin}[1]{\texttt{#1b}}
\newcommand{\dec}[1]{\texttt{#1}}

\newcommand{\important}[1]{\textbf{Penting:} #1}
\newcommand{\catatan}[1]{\textbf{Catatan:} #1}
\newcommand{\term}[1]{\textbf{#1}}

\newcommand{\CPL}[1]{\textbf{CPL#1}}
\newcommand{\CPMK}[1]{\textbf{CPMK#1}}
\newcommand{\SubCPMK}[1]{\textbf{Sub-CPMK #1}}

\newcommand{\checkbox}{\ding{113}\hspace{0.35em}}
\newcommand{\checked}{\ding{51}\hspace{0.35em}}
\newcommand{\crossed}{\ding{55}\hspace{0.35em}}

\newcommand{\flowarrow}{\ensuremath{\rightarrow}}

% ---------------------------------------------------------
% 6. Makro tata letak slide
% ---------------------------------------------------------

% Keterangan sumber materi (ketertelusuran ke buku)
\newcommand{\Sumber}[1]{%
    \vfill
    {\scriptsize\color{obeGray}Sumber: #1\par}%
}

% Frame judul deck: \JudulDeck{<judul bab>}{<pertemuan>}
\newcommand{\JudulDeck}[2]{%
\begin{frame}[plain]
    \vspace{0.25cm}
    {\color{obeBlue}\sffamily\bfseries\small \NamaMK\ (\KodeMK)\par}
    \vspace{0.12cm}
    {\color{obeGray}\sffamily\scriptsize
     \Prodi\ \textbullet\ \Fakultas\ \textbullet\ \PerguruanTinggi\par}
    \vfill
    {\color{obeBlue}\sffamily\LARGE\bfseries #1\par}
    \vspace{0.4cm}
    {\color{obeTeal}\sffamily\large #2\par}
    \vspace{0.45cm}
    {\color{obeTeal}\rule{\textwidth}{1pt}}
    \vfill
    {\color{obeGray}\sffamily\small Dosen Pengampu\par}
    \vspace{0.08cm}
    {\sffamily\bfseries \DosenPengampu\par}
    \vspace{0.2cm}
    {\color{obeGray}\sffamily\small \SemesterTeks\par}
    \vspace{0.25cm}
\end{frame}
}

% Frame pemisah bagian
\newcommand{\PemisahBagian}[1]{%
\begin{frame}[plain]
    \vfill
    \begin{center}
        {\color{obeTeal}\sffamily\small\bfseries BAGIAN\par}
        \vspace{0.35cm}
        {\color{obeBlue}\sffamily\Large\bfseries #1\par}
        \vspace{0.35cm}
        {\color{obeTeal}\rule{0.55\textwidth}{0.9pt}}
    \end{center}
    \vfill
\end{frame}
}

% ---------------------------------------------------------
% 7. Gaya penampil kode
% ---------------------------------------------------------
% =========================================================
% ppt/src/listing-ppt.tex
% Gaya penampil kode untuk slide.
%
% Diturunkan dari buku_ajar/config/listings.tex, tetapi
% disesuaikan untuk slide: ukuran lebih kecil (karena frame
% 16:9 dan kode harus utuh terlihat) dan jarak lebih rapat.
% =========================================================

% ---------------------------------------------------------
% 1. Definisi bahasa MASM/TASM 8086 (salinan dari buku)
% ---------------------------------------------------------
\lstdefinelanguage{MASM}{
    morekeywords={
        MOV, XCHG, LEA, LDS, LES, PUSH, POP, PUSHA, POPA, PUSHF, POPF,
        LAHF, SAHF, XLAT, IN, OUT, CBW, CWD,
        ADD, ADC, SUB, SBB, INC, DEC, NEG, CMP, MUL, IMUL, DIV, IDIV,
        AAA, AAS, AAM, AAD, DAA, DAS,
        AND, OR, XOR, NOT, TEST, SHL, SHR, SAL, SAR, ROL, ROR, RCL, RCR,
        JMP, JE, JZ, JNE, JNZ, JG, JGE, JL, JLE, JA, JAE, JB, JBE,
        JC, JNC, JS, JNS, JO, JNO, JP, JPE, JNP, JPO, JCXZ,
        LOOP, LOOPE, LOOPZ, LOOPNE, LOOPNZ, CALL, RET, INT, IRET, INTO,
        HLT, NOP, CLI, STI, CLD, STD, CLC, STC, CMC, WAIT, ESC, LOCK,
        DB, DW, DD, DQ, DT, EQU, ORG, SEGMENT, ENDS, ASSUME, END,
        PROC, ENDP, MACRO, ENDM, LOCAL, PUBLIC, EXTRN, EXTERN,
        INCLUDE, INCLUDELIB, MODEL, STACK, CODE, DATA, CONST, ALIGN, EVEN,
        OFFSET, PTR, BYTE, WORD, DWORD, QWORD, TBYTE,
        NEAR, FAR, SHORT, LABEL, TYPE, SIZE, LENGTH, DUP, THIS,
        AX, BX, CX, DX, AL, AH, BL, BH, CL, CH, DL, DH,
        SP, BP, SI, DI, IP, CS, DS, SS, ES, FLAGS,
        CF, PF, AF, ZF, TF, IF, DF, OF
    },
    sensitive=false,
    morecomment=[l]{;},
    morestring=[b]',
    morestring=[b]"
}

% ---------------------------------------------------------
% 2. Gaya kode assembly untuk slide
% ---------------------------------------------------------
\lstdefinestyle{asmslide}{
    language=MASM,
    basicstyle=\ttfamily\scriptsize,
    keywordstyle=\color{obeBlue}\bfseries,
    commentstyle=\color{obeGreen}\itshape,
    stringstyle=\color{obeMaroon},
    numberstyle=\tiny\color{obeGray},
    numbers=left,
    numbersep=5pt,
    stepnumber=1,
    backgroundcolor=\color{codeBack},
    frame=single,
    rulecolor=\color{codeFrame},
    framesep=4pt,
    breaklines=true,
    breakatwhitespace=false,
    % CATATAN PENTING: opsi `postbreak` TIDAK boleh dipakai di sini.
    % Di beamer, `breaklines=true` bersama `postbreak=<kode>` membuat
    % listings gagal dirakit dengan galat
    %   "Missing number, treated as zero <to be read again> \kern".
    % Karena \lstset{style=asmslide} bersifat global, opsi postbreak
    % juga bocor ke gaya lain (mis. outslide/keluaran) dan merusaknya.
    % Di buku ajar (kelas `book`) justru sebaliknya: postbreak aman dan
    % dipakai. Jadi larangan ini khusus untuk slide.
    showstringspaces=false,
    tabsize=4,
    columns=fullflexible,
    keepspaces=true,
    upquote=true,
    xleftmargin=12pt,
    aboveskip=0.4em,
    belowskip=0.4em
}

% ---------------------------------------------------------
% 3. Gaya keluaran program / sesi debugger untuk slide
% ---------------------------------------------------------
\lstdefinestyle{outslide}{
    basicstyle=\ttfamily\scriptsize,
    keywordstyle=\color{black},
    commentstyle=\color{obeGray},
    stringstyle=\color{black},
    numbers=none,
    backgroundcolor=\color{obeGrayLight},
    frame=single,
    rulecolor=\color{codeFrame},
    framesep=4pt,
    breaklines=true,
    showstringspaces=false,
    columns=fullflexible,
    keepspaces=true,
    upquote=true,
    xleftmargin=4pt,
    aboveskip=0.4em,
    belowskip=0.4em
}

\lstset{style=asmslide}

% ---------------------------------------------------------
% 4. Environment ringkas untuk dipakai di dalam frame
%
% Frame yang memuat kedua environment ini WAJIB memakai
% opsi [fragile], karena isinya verbatim.
% ---------------------------------------------------------
\lstnewenvironment{kodeasm}[1][]{%
    \lstset{style=asmslide, #1}%
}{}

\lstnewenvironment{keluaran}[1][]{%
    \lstset{style=outslide, #1}%
}{}


%
% Palet warna disalin dari buku_ajar/config/styles.tex agar
% identitas visual slide sama dengan buku ajar.
% =========================================================

% ---------------------------------------------------------
% 1. Paket
% ---------------------------------------------------------
\usepackage[T1]{fontenc}
\usepackage[utf8]{inputenc}
\usepackage[indonesian]{babel}
\usepackage{lmodern}
\usepackage{tikz}
\usetikzlibrary{arrows.meta, positioning, calc, shapes.geometric}
\usepackage{booktabs}
\usepackage{array}
\usepackage{amsmath}
\usepackage{amsfonts}
\usepackage{pifont}          % \ding untuk tanda checklist
\usepackage{listings}

% ---------------------------------------------------------
% 2. Palet warna OBE (identik dengan buku)
% ---------------------------------------------------------
\definecolor{obeBlue}{HTML}{1F3864}
\definecolor{obeBlueLight}{HTML}{DCE4F5}
\definecolor{obeTeal}{HTML}{0F6E6E}
\definecolor{obeTealLight}{HTML}{D9F0EE}
\definecolor{obeGreen}{HTML}{2E6B34}
\definecolor{obeGreenLight}{HTML}{DEEFDF}
\definecolor{obeAmber}{HTML}{8A5A00}
\definecolor{obeAmberLight}{HTML}{FDF1D6}
\definecolor{obeMaroon}{HTML}{7A2233}
\definecolor{obeMaroonLight}{HTML}{F7E1E5}
\definecolor{obeGray}{HTML}{4A4A4A}
\definecolor{obeGrayLight}{HTML}{F0F0F0}
\definecolor{codeBack}{HTML}{F7F7F4}
\definecolor{codeFrame}{HTML}{C9C9C0}

% ---------------------------------------------------------
% 3. Identitas mata kuliah (dipakai seluruh deck)
% ---------------------------------------------------------
\newcommand{\NamaMK}{Bahasa Pemrograman Tingkat Rendah}
\newcommand{\NamaMKShort}{Bhs.\ Pemrograman Tingkat Rendah}
\newcommand{\KodeMK}{MWCA010}
\newcommand{\Prodi}{Program Studi Teknik Informatika}
\newcommand{\Fakultas}{Fakultas Teknologi Informasi}
\newcommand{\PerguruanTinggi}{Universitas Bale Bandung}
\newcommand{\DosenPengampu}{Cecep Suwanda, S.Si., M.Kom.}
\newcommand{\SemesterTeks}{Semester Ganjil}

% \PertemuanTeks dan \JudulBab ditetapkan di masing-masing
% berkas minggu-NN.tex sebelum \begin{document}.
\providecommand{\PertemuanTeks}{Pertemuan}

% ---------------------------------------------------------
% 4. Tema beamer
% ---------------------------------------------------------
\usetheme{default}
\setbeamertemplate{navigation symbols}{}

\setbeamercolor{normal text}{fg=black!88}
\setbeamercolor{structure}{fg=obeBlue}
\setbeamercolor{alerted text}{fg=obeMaroon}

% --- judul slide: pita biru penuh lebar ---
\setbeamercolor{frametitle}{fg=white,bg=obeBlue}
\setbeamerfont{frametitle}{series=\bfseries,size=\large}
\setbeamertemplate{frametitle}{%
    \nointerlineskip%
    \begin{beamercolorbox}[%
        wd=\paperwidth,ht=2.9ex,dp=1.5ex,
        leftskip=0.4cm,rightskip=0.4cm]{frametitle}%
        \usebeamerfont{frametitle}\insertframetitle%
    \end{beamercolorbox}%
    \nointerlineskip%
    \hbox{\textcolor{obeTeal}{\rule{\paperwidth}{1.1pt}}}%
}

% --- kaki slide: identitas mata kuliah, pertemuan, nomor ---
\setbeamercolor{footline}{fg=white,bg=obeGray}
\setbeamerfont{footline}{size=\scriptsize,family=\sffamily}
\setbeamertemplate{footline}{%
    \leavevmode%
    \hbox{%
        \begin{beamercolorbox}[wd=0.42\paperwidth,ht=2.1ex,dp=1.1ex,
            leftskip=0.4cm]{footline}%
            \usebeamerfont{footline}%
            \KodeMK\ \textbullet\ \NamaMKShort%
        \end{beamercolorbox}%
        \begin{beamercolorbox}[wd=0.32\paperwidth,ht=2.1ex,dp=1.1ex,
            center]{footline}%
            \usebeamerfont{footline}\PertemuanTeks%
        \end{beamercolorbox}%
        \begin{beamercolorbox}[wd=0.26\paperwidth,ht=2.1ex,dp=1.1ex,
            rightskip=0.4cm]{footline}%
            \usebeamerfont{footline}%
            \insertframenumber\,/\,\inserttotalframenumber%
        \end{beamercolorbox}%
    }%
    \vskip0pt%
}

% --- daftar ---
\setbeamercolor{itemize item}{fg=obeTeal}
\setbeamercolor{itemize subitem}{fg=obeTeal}
\setbeamercolor{enumerate item}{fg=obeBlue}
\setbeamertemplate{itemize item}{\raise0.2ex\hbox{\tiny$\blacksquare$}}
\setbeamertemplate{itemize subitem}{\raise0.2ex\hbox{\scriptsize$\bullet$}}

% --- kotak (pengganti tcolorbox pada buku) ---
\setbeamercolor{block title}{bg=obeBlueLight,fg=obeBlue}
\setbeamercolor{block body}{bg=obeBlueLight!45,fg=black!88}
\setbeamercolor{block title alerted}{bg=obeMaroonLight,fg=obeMaroon}
\setbeamercolor{block body alerted}{bg=obeMaroonLight!45,fg=black!88}
\setbeamercolor{block title example}{bg=obeTealLight,fg=obeTeal}
\setbeamercolor{block body example}{bg=obeTealLight!45,fg=black!88}

% ---------------------------------------------------------
% 5. Makro notasi (sama persis dengan buku)
%
% Catatan: \note TIDAK didefinisikan di sini karena beamer
% sudah memakai \note untuk catatan pembicara. Untuk catatan
% di slide dipakai \catatan.
% ---------------------------------------------------------
\newcommand{\code}[1]{\texttt{#1}}
\newcommand{\asm}[1]{\texttt{#1}}
\newcommand{\reg}[1]{\texttt{#1}}
\newcommand{\flag}[1]{\texttt{#1}}
\newcommand{\dir}[1]{\texttt{#1}}
\newcommand{\file}[1]{\texttt{#1}}
\newcommand{\eng}[1]{\textit{#1}}

\newcommand{\hex}[1]{\texttt{#1h}}
\newcommand{\bin}[1]{\texttt{#1b}}
\newcommand{\dec}[1]{\texttt{#1}}

\newcommand{\important}[1]{\textbf{Penting:} #1}
\newcommand{\catatan}[1]{\textbf{Catatan:} #1}
\newcommand{\term}[1]{\textbf{#1}}

\newcommand{\CPL}[1]{\textbf{CPL#1}}
\newcommand{\CPMK}[1]{\textbf{CPMK#1}}
\newcommand{\SubCPMK}[1]{\textbf{Sub-CPMK #1}}

\newcommand{\checkbox}{\ding{113}\hspace{0.35em}}
\newcommand{\checked}{\ding{51}\hspace{0.35em}}
\newcommand{\crossed}{\ding{55}\hspace{0.35em}}

\newcommand{\flowarrow}{\ensuremath{\rightarrow}}

% ---------------------------------------------------------
% 6. Makro tata letak slide
% ---------------------------------------------------------

% Keterangan sumber materi (ketertelusuran ke buku)
\newcommand{\Sumber}[1]{%
    \vfill
    {\scriptsize\color{obeGray}Sumber: #1\par}%
}

% Frame judul deck: \JudulDeck{<judul bab>}{<pertemuan>}
\newcommand{\JudulDeck}[2]{%
\begin{frame}[plain]
    \vspace{0.25cm}
    {\color{obeBlue}\sffamily\bfseries\small \NamaMK\ (\KodeMK)\par}
    \vspace{0.12cm}
    {\color{obeGray}\sffamily\scriptsize
     \Prodi\ \textbullet\ \Fakultas\ \textbullet\ \PerguruanTinggi\par}
    \vfill
    {\color{obeBlue}\sffamily\LARGE\bfseries #1\par}
    \vspace{0.4cm}
    {\color{obeTeal}\sffamily\large #2\par}
    \vspace{0.45cm}
    {\color{obeTeal}\rule{\textwidth}{1pt}}
    \vfill
    {\color{obeGray}\sffamily\small Dosen Pengampu\par}
    \vspace{0.08cm}
    {\sffamily\bfseries \DosenPengampu\par}
    \vspace{0.2cm}
    {\color{obeGray}\sffamily\small \SemesterTeks\par}
    \vspace{0.25cm}
\end{frame}
}

% Frame pemisah bagian
\newcommand{\PemisahBagian}[1]{%
\begin{frame}[plain]
    \vfill
    \begin{center}
        {\color{obeTeal}\sffamily\small\bfseries BAGIAN\par}
        \vspace{0.35cm}
        {\color{obeBlue}\sffamily\Large\bfseries #1\par}
        \vspace{0.35cm}
        {\color{obeTeal}\rule{0.55\textwidth}{0.9pt}}
    \end{center}
    \vfill
\end{frame}
}

% ---------------------------------------------------------
% 7. Gaya penampil kode
% ---------------------------------------------------------
% =========================================================
% ppt/src/listing-ppt.tex
% Gaya penampil kode untuk slide.
%
% Diturunkan dari buku_ajar/config/listings.tex, tetapi
% disesuaikan untuk slide: ukuran lebih kecil (karena frame
% 16:9 dan kode harus utuh terlihat) dan jarak lebih rapat.
% =========================================================

% ---------------------------------------------------------
% 1. Definisi bahasa MASM/TASM 8086 (salinan dari buku)
% ---------------------------------------------------------
\lstdefinelanguage{MASM}{
    morekeywords={
        MOV, XCHG, LEA, LDS, LES, PUSH, POP, PUSHA, POPA, PUSHF, POPF,
        LAHF, SAHF, XLAT, IN, OUT, CBW, CWD,
        ADD, ADC, SUB, SBB, INC, DEC, NEG, CMP, MUL, IMUL, DIV, IDIV,
        AAA, AAS, AAM, AAD, DAA, DAS,
        AND, OR, XOR, NOT, TEST, SHL, SHR, SAL, SAR, ROL, ROR, RCL, RCR,
        JMP, JE, JZ, JNE, JNZ, JG, JGE, JL, JLE, JA, JAE, JB, JBE,
        JC, JNC, JS, JNS, JO, JNO, JP, JPE, JNP, JPO, JCXZ,
        LOOP, LOOPE, LOOPZ, LOOPNE, LOOPNZ, CALL, RET, INT, IRET, INTO,
        HLT, NOP, CLI, STI, CLD, STD, CLC, STC, CMC, WAIT, ESC, LOCK,
        DB, DW, DD, DQ, DT, EQU, ORG, SEGMENT, ENDS, ASSUME, END,
        PROC, ENDP, MACRO, ENDM, LOCAL, PUBLIC, EXTRN, EXTERN,
        INCLUDE, INCLUDELIB, MODEL, STACK, CODE, DATA, CONST, ALIGN, EVEN,
        OFFSET, PTR, BYTE, WORD, DWORD, QWORD, TBYTE,
        NEAR, FAR, SHORT, LABEL, TYPE, SIZE, LENGTH, DUP, THIS,
        AX, BX, CX, DX, AL, AH, BL, BH, CL, CH, DL, DH,
        SP, BP, SI, DI, IP, CS, DS, SS, ES, FLAGS,
        CF, PF, AF, ZF, TF, IF, DF, OF
    },
    sensitive=false,
    morecomment=[l]{;},
    morestring=[b]',
    morestring=[b]"
}

% ---------------------------------------------------------
% 2. Gaya kode assembly untuk slide
% ---------------------------------------------------------
\lstdefinestyle{asmslide}{
    language=MASM,
    basicstyle=\ttfamily\scriptsize,
    keywordstyle=\color{obeBlue}\bfseries,
    commentstyle=\color{obeGreen}\itshape,
    stringstyle=\color{obeMaroon},
    numberstyle=\tiny\color{obeGray},
    numbers=left,
    numbersep=5pt,
    stepnumber=1,
    backgroundcolor=\color{codeBack},
    frame=single,
    rulecolor=\color{codeFrame},
    framesep=4pt,
    breaklines=true,
    breakatwhitespace=false,
    % CATATAN PENTING: opsi `postbreak` TIDAK boleh dipakai di sini.
    % Di beamer, `breaklines=true` bersama `postbreak=<kode>` membuat
    % listings gagal dirakit dengan galat
    %   "Missing number, treated as zero <to be read again> \kern".
    % Karena \lstset{style=asmslide} bersifat global, opsi postbreak
    % juga bocor ke gaya lain (mis. outslide/keluaran) dan merusaknya.
    % Di buku ajar (kelas `book`) justru sebaliknya: postbreak aman dan
    % dipakai. Jadi larangan ini khusus untuk slide.
    showstringspaces=false,
    tabsize=4,
    columns=fullflexible,
    keepspaces=true,
    upquote=true,
    xleftmargin=12pt,
    aboveskip=0.4em,
    belowskip=0.4em
}

% ---------------------------------------------------------
% 3. Gaya keluaran program / sesi debugger untuk slide
% ---------------------------------------------------------
\lstdefinestyle{outslide}{
    basicstyle=\ttfamily\scriptsize,
    keywordstyle=\color{black},
    commentstyle=\color{obeGray},
    stringstyle=\color{black},
    numbers=none,
    backgroundcolor=\color{obeGrayLight},
    frame=single,
    rulecolor=\color{codeFrame},
    framesep=4pt,
    breaklines=true,
    showstringspaces=false,
    columns=fullflexible,
    keepspaces=true,
    upquote=true,
    xleftmargin=4pt,
    aboveskip=0.4em,
    belowskip=0.4em
}

\lstset{style=asmslide}

% ---------------------------------------------------------
% 4. Environment ringkas untuk dipakai di dalam frame
%
% Frame yang memuat kedua environment ini WAJIB memakai
% opsi [fragile], karena isinya verbatim.
% ---------------------------------------------------------
\lstnewenvironment{kodeasm}[1][]{%
    \lstset{style=asmslide, #1}%
}{}

\lstnewenvironment{keluaran}[1][]{%
    \lstset{style=outslide, #1}%
}{}



\renewcommand{\PertemuanTeks}{Pertemuan 13}

\begin{document}

% =========================================================
% 1. JUDUL
% =========================================================
\JudulDeck{Stack, Prosedur, dan Macro}{Pertemuan 13}

% =========================================================
% 2. CAPAIAN DAN TUJUAN
% =========================================================
\begin{frame}{Capaian dan Tujuan Pembelajaran}
\begin{block}{Sub-CPMK 3.5}
\footnotesize Menggunakan stack (PUSH/POP), prosedur
(PROC/ENDP, CALL/RET), dan macro dalam program yang modular.
\end{block}

\vspace{0.1em}
{\small\textbf{CPMK yang didukung:} \CPMK{03}\quad\textbullet\quad
\textbf{CPL yang didukung:} \CPL{03}\par}

\vspace{0.15em}
\begin{exampleblock}{Setelah pertemuan ini, mahasiswa mampu:}
\footnotesize
\begin{enumerate}
    \item menjelaskan sifat \eng{last in, first out} pada stack;
    \item menggunakan PUSH/POP dan menjelaskan perubahan SP;
    \item menulis prosedur dengan PROC/ENDP serta CALL/RET;
    \item mengakses parameter melalui \eng{stack frame} dan BP;
    \item mendefinisikan macro berparameter dan membedakannya dari
          prosedur.
\end{enumerate}
\end{exampleblock}
\end{frame}

% =========================================================
% 3. PETA MATERI
% =========================================================
\begin{frame}{Peta Materi}
\small
\begin{enumerate}
    \item \textbf{Stack dan Instruksi PUSH/POP}\\
          {\footnotesize LIFO; SS:SP; PUSH dan POP; aturan
          pemakaian}

    \vspace{0.25em}
    \item \textbf{Prosedur dan Penyampaian Data}\\
          {\footnotesize PROC/ENDP; CALL/RET; stack frame dan BP;
          register/stack/variabel}

    \vspace{0.25em}
    \item \textbf{Macro}\\
          {\footnotesize MACRO/ENDM; parameter; perbedaan dengan
          prosedur}
\end{enumerate}

\vspace{0.35em}
\begin{block}{Praktik pada pertemuan ini}
\small Menelusuri stack, prosedur berparameter, dan macro.
\end{block}
\end{frame}

% =========================================================
% BAGIAN 1
% =========================================================
\PemisahBagian{Stack dan Instruksi PUSH/POP}

% ---------------------------------------------------------
\begin{frame}{Stack: LIFO dan SS:SP}
\small
\begin{itemize}
    \item Stack menyimpan data sementara dengan aturan
          \eng{last in, first out}.

    \vspace{0.18em}
    \item Stack dipakai untuk alamat kembali, parameter, dan nilai
          sementara.

    \vspace{0.18em}
    \item Puncak stack ditunjuk \reg{SS}:\reg{SP}; alamat fisiknya
          $(\reg{SS} \times 16) + \reg{SP}$.

    \vspace{0.18em}
    \item Stack 8086 tumbuh ke alamat yang lebih rendah; satu sel
          berukuran satu word.
\end{itemize}

\vspace{0.2em}
\begin{block}{Arah pertumbuhan}
\small Setiap data didorong, SP berkurang dua; setiap data diambil,
SP bertambah dua.
\end{block}
\Sumber{Bab XIV, Bagian 1}
\end{frame}

% ---------------------------------------------------------
\begin{frame}[fragile]{PUSH}
\small
\code{PUSH} menyimpan satu word ke puncak stack:
SP dikurangi dua, lalu operand disalin ke \reg{SS}:\reg{SP}.

\vspace{0.2em}
\begin{kodeasm}[numbers=none]
MOV AX, 1234h
PUSH AX             ; SP = SP - 2; [SS:SP] = 1234h
\end{kodeasm}

\vspace{0.15em}
{\footnotesize\color{obeGray}PUSH hanya bekerja pada operand 16 bit;
pada 8086 tidak ada \code{PUSH AL} atau nilai segera.\par}
\Sumber{Bab XIV, Bagian 1}
\end{frame}

% ---------------------------------------------------------
\begin{frame}[fragile]{POP}
\small
\code{POP} mengambil satu word dari puncak stack:
isi \reg{SS}:\reg{SP} disalin ke tujuan, lalu SP ditambah dua.

\vspace{0.2em}
\begin{kodeasm}[numbers=none]
POP BX              ; BX = [SS:SP]; SP = SP + 2
\end{kodeasm}

\vspace{0.15em}
{\footnotesize\color{obeGray}Setiap PUSH harus berpasangan dengan
POP agar SP kembali ke posisi semula.\par}
\Sumber{Bab XIV, Bagian 1}
\end{frame}

% ---------------------------------------------------------
\begin{frame}{Perubahan Stack}
\begin{center}
\begin{tikzpicture}[
    sel/.style={
        draw=obeBlue, line width=0.6pt, minimum width=3.4cm,
        minimum height=0.55cm, font=\ttfamily\small, inner sep=0pt
    },
    isi/.style={sel, fill=obeBlueLight},
    kosong/.style={sel, fill=obeGrayLight},
    panah/.style={-{Stealth[length=2.2mm]}, line width=0.9pt,
        obeMaroon}
]
    \node[isi] (a0) at (0,0) {1234h};
    \node[kosong] (a1) at (0,-0.55) {kosong};
    \node[font=\scriptsize, text=obeGray, anchor=west]
        at (1.9,0) {SP sesudah PUSH};
    \node[font=\scriptsize, text=obeGray, anchor=west]
        at (1.9,-0.55) {SP sebelum PUSH};

    \draw[panah] (-2.1,0) -- (-1.8,0);
    \draw[panah] (-2.1,-0.55) -- (-1.8,-0.55);
    \node[font=\scriptsize, text=obeMaroon, anchor=east]
        at (-2.2,0) {PUSH};
    \node[font=\scriptsize, text=obeMaroon, anchor=east]
        at (-2.2,-0.55) {awal};
\end{tikzpicture}
\end{center}

\vspace{0.2em}
{\footnotesize\color{obeGray}POP bergerak sebaliknya: data teratas
diambil lebih dahulu, lalu SP bertambah dua.\par}
\Sumber{Bab XIV, Bagian 1}
\end{frame}

% ---------------------------------------------------------
\begin{frame}{Aturan Pemakaian Stack}
\small
\begin{itemize}
    \item PUSH dan POP bekerja pada operand 16 bit.

    \vspace{0.18em}
    \item Nilai segera tidak dapat didorong langsung pada 8086.

    \vspace{0.18em}
    \item Jumlah PUSH dan POP harus seimbang.

    \vspace{0.18em}
    \item Isi stack bersifat sementara dan dapat ditimpa pemanggilan
          berikutnya.
\end{itemize}

\vspace{0.25em}
\begin{alertblock}{Bila tidak seimbang}
\small SP tidak kembali ke posisi semula; RET dapat mengambil alamat
kembali yang salah.
\end{alertblock}
\Sumber{Bab XIV, Bagian 1}
\end{frame}

% ---------------------------------------------------------
\begin{frame}{Rangkuman Bagian 1}
\small
\begin{itemize}
    \item Stack bersifat LIFO dan ditunjuk SS:SP.

    \vspace{0.15em}
    \item Stack tumbuh ke alamat rendah; satu sel dua byte.

    \vspace{0.15em}
    \item PUSH: SP $-2$, lalu simpan operand.

    \vspace{0.15em}
    \item POP: ambil operand, lalu SP $+2$.

    \vspace{0.15em}
    \item Jumlah PUSH dan POP harus seimbang.
\end{itemize}
\end{frame}

% =========================================================
% BAGIAN 2
% =========================================================
\PemisahBagian{Prosedur dan Penyampaian Data}

% ---------------------------------------------------------
\begin{frame}[fragile]{PROC, ENDP, CALL, dan RET}
\small
Prosedur didefinisikan dengan PROC dan diakhiri ENDP. Pemanggil
memakai CALL; prosedur kembali dengan RET.

\vspace{0.2em}
\begin{kodeasm}[numbers=none]
HitungLuas PROC
    ; isi prosedur
    RET
HitungLuas ENDP

CALL HitungLuas
\end{kodeasm}

\vspace{0.1em}
{\footnotesize\color{obeGray}CALL menyimpan alamat kembali; RET
memuatnya kembali ke IP.\par}
\Sumber{Bab XIV, Bagian 2}
\end{frame}

% ---------------------------------------------------------
\begin{frame}{CALL vs JMP}
\footnotesize
CALL menyimpan alamat instruksi sesudahnya ke stack; JMP tidak.

\vspace{0.2em}
\begin{center}
\begin{tabular}{@{}>{\raggedright\arraybackslash}p{3.0cm}
    >{\raggedright\arraybackslash}p{5.0cm}
    >{\raggedright\arraybackslash}p{4.8cm}@{}}
    \toprule
    \textbf{Aspek} & \textbf{CALL dekat} & \textbf{CALL jauh} \\
    \midrule
    Alamat disimpan & IP (2 byte) & CS dan IP (4 byte) \\
    \addlinespace
    Instruksi kembali & RET & RETF \\
    \addlinespace
    Pemakaian & Segmen kode sama & Segmen kode lain \\
    \bottomrule
\end{tabular}
\end{center}

\vspace{0.2em}
{\color{obeGray}Contoh pada bab ini memakai pemanggilan dekat.\par}
\Sumber{Bab XIV, Bagian 2}
\end{frame}

% ---------------------------------------------------------
\begin{frame}{Stack Frame dan Register BP}
\begin{center}
\begin{tikzpicture}[
    sel/.style={
        draw=obeBlue, line width=0.7pt, minimum width=4.6cm,
        minimum height=0.55cm, align=center, font=\small,
        inner sep=2pt
    },
    param/.style={sel, fill=obeAmberLight, draw=obeAmber},
    kendali/.style={sel, fill=obeTealLight, draw=obeTeal},
    lama/.style={sel, fill=obeGrayLight, draw=obeGray},
    ket/.style={font=\scriptsize, text=obeGray, anchor=west}
]
    \node[param] (p2) at (0,0) {parameter kedua \quad [BP+6]};
    \node[param] (p1) at (0,-0.55) {parameter pertama \quad [BP+4]};
    \node[kendali] (ret) at (0,-1.1) {alamat kembali \quad [BP+2]};
    \node[lama] (bp) at (0,-1.65) {BP lama \quad [BP+0]};
    \node[sel, fill=obeBlueLight] (lokal) at (0,-2.2)
        {variabel lokal / ruang bebas};

    \node[ket] at (2.5,0) {alamat tinggi};
    \node[ket] at (2.5,-2.2) {alamat rendah};
    \node[font=\scriptsize, text=obeMaroon, anchor=east]
        at (-2.5,-1.65) {BP dan SP};
    \draw[-{Stealth[length=2.2mm]}, line width=0.9pt, obeMaroon]
        (-2.5,-1.65) -- (-2.35,-1.65);
\end{tikzpicture}
\end{center}

\vspace{0.1em}
{\footnotesize\color{obeGray}BP dibentuk dengan PUSH BP lalu MOV BP,
SP.\par}
\Sumber{Bab XIV, Bagian 2}
\end{frame}

% ---------------------------------------------------------
\begin{frame}{Akses Parameter melalui BP}
\footnotesize
\begin{center}
\begin{tabular}{@{}>{\raggedright\arraybackslash}p{3.0cm}
    >{\raggedright\arraybackslash}p{9.0cm}@{}}
    \toprule
    \textbf{Alamat} & \textbf{Isi} \\
    \midrule
    \code{[BP+0]} & Nilai BP lama \\
    \addlinespace
    \code{[BP+2]} & Alamat kembali \\
    \addlinespace
    \code{[BP+4]} & Parameter pertama \\
    \addlinespace
    \code{[BP+6]} & Parameter kedua \\
    \bottomrule
\end{tabular}
\end{center}

\vspace{0.2em}
{\color{obeGray}Parameter pertama berada paling dekat dengan alamat
kembali karena didorong terakhir.\par}
\Sumber{Bab XIV, Bagian 2}
\end{frame}

% ---------------------------------------------------------
\begin{frame}{Penyampaian Data}
\footnotesize
\begin{center}
\begin{tabular}{@{}>{\raggedright\arraybackslash}p{2.4cm}
    >{\raggedright\arraybackslash}p{4.4cm}
    >{\raggedright\arraybackslash}p{5.8cm}@{}}
    \toprule
    \textbf{Cara} & \textbf{Contoh} & \textbf{Keterangan} \\
    \midrule
    Register & \code{MOV AX, lebar}; \code{CALL Hitung} &
    Cepat; jumlah parameter terbatas \\
    \addlinespace
    Stack & \code{PUSH lebar}; \code{CALL Hitung} &
    Banyak parameter; akses \code{[BP+4]} \\
    \addlinespace
    Variabel & \code{MOV lebar, 5}; \code{CALL Hitung} &
    Sederhana; bergantung variabel global \\
    \bottomrule
\end{tabular}
\end{center}

\vspace{0.2em}
{\color{obeGray}RET $n$ membuang $n$ byte parameter setelah
kembali.\par}
\Sumber{Bab XIV, Bagian 2}
\end{frame}

% ---------------------------------------------------------
\begin{frame}{Rangkuman Bagian 2}
\small
\begin{itemize}
    \item Prosedur memakai PROC/ENDP; CALL memanggil, RET kembali.

    \vspace{0.15em}
    \item CALL menyimpan alamat kembali ke stack.

    \vspace{0.15em}
    \item Stack frame dibentuk PUSH BP lalu MOV BP, SP.

    \vspace{0.15em}
    \item Parameter pertama di [BP+4], kedua di [BP+6].

    \vspace{0.15em}
    \item Data dapat dikirim melalui register, stack, atau variabel.
\end{itemize}
\end{frame}

% =========================================================
% BAGIAN 3
% =========================================================
\PemisahBagian{Macro}

% ---------------------------------------------------------
\begin{frame}[fragile]{Mendefinisikan Macro}
\small
Macro didefinisikan dengan MACRO dan diakhiri ENDM. Assembler
menyalin isi macro ke tempat pemakaian saat perakitan.

\vspace{0.2em}
\begin{kodeasm}[numbers=none]
Cetak_Kar MACRO Kar
    MOV CX, 3
    MOV AH, 02h
    MOV DL, Kar
Ulang:
    INT 21h
    LOOP Ulang
ENDM

Cetak_Kar 'S'
\end{kodeasm}
\Sumber{Bab XIV, Bagian 3}
\end{frame}

% ---------------------------------------------------------
\begin{frame}{Parameter Macro}
\small
\begin{itemize}
    \item Parameter macro diganti secara tekstual oleh assembler.

    \vspace{0.18em}
    \item Argumen dapat berupa karakter, angka, nama variabel, atau
          register.

    \vspace{0.18em}
    \item Prosedur tidak memiliki daftar parameter pada PROC; data
          dikirim lewat register, stack, atau variabel.

    \vspace{0.18em}
    \item Macro memiliki daftar parameter langsung pada definisinya.
\end{itemize}

\vspace{0.2em}
\begin{alertblock}{Label dalam macro}
\small Bila macro berlabel dipakai berkali-kali, label yang sama
disalin berulang dan dapat memicu \file{Symbol already defined}.
\end{alertblock}
\Sumber{Bab XIV, Bagian 3}
\end{frame}

% ---------------------------------------------------------
\begin{frame}{Prosedur vs Macro}
\footnotesize
\begin{center}
\begin{tabular}{@{}>{\raggedright\arraybackslash}p{3.0cm}
    >{\raggedright\arraybackslash}p{5.0cm}
    >{\raggedright\arraybackslash}p{4.8cm}@{}}
    \toprule
    \textbf{Aspek} & \textbf{Prosedur} & \textbf{Macro} \\
    \midrule
    Pemanggilan & CALL dan RET & nama macro ditulis langsung \\
    \addlinespace
    Ukuran program & kode ditulis sekali &
    kode disalin setiap pemakaian \\
    \addlinespace
    Kecepatan & ada biaya CALL/RET & tanpa biaya pemanggilan \\
    \addlinespace
    Parameter & register, stack, variabel &
    daftar parameter pada definisi \\
    \addlinespace
    Segmen & terikat segmen untuk call dekat &
    mengikuti tempat pemakaian \\
    \bottomrule
\end{tabular}
\end{center}
\Sumber{Bab XIV, Bagian 3}
\end{frame}

% ---------------------------------------------------------
\begin{frame}{Kapan Memilih?}
\small
\begin{itemize}
    \item Pilih prosedur bila program menekankan ukuran dan
          modularitas.

    \vspace{0.18em}
    \item Pilih macro bila bagian tersebut pendek, sering dipakai,
          dan parameternya membuat penulisan lebih ringkas.

    \vspace{0.18em}
    \item Macro memperbesar ukuran program karena isinya disalin
          setiap kali dipakai.

    \vspace{0.18em}
    \item Prosedur lebih tepat untuk program besar dengan banyak
          pemanggilan.
\end{itemize}

\vspace{0.2em}
\begin{block}{Ringkas}
\small Prosedur untuk ukuran; macro untuk keringkasan penulisan.
\end{block}
\Sumber{Bab XIV, Bagian 3}
\end{frame}

% ---------------------------------------------------------
\begin{frame}{Rangkuman Bagian 3}
\small
\begin{itemize}
    \item Macro didefinisikan dengan MACRO dan diakhiri ENDM.

    \vspace{0.15em}
    \item Parameter macro diganti secara tekstual saat perakitan.

    \vspace{0.15em}
    \item Macro disalin ke setiap tempat pemakaian.

    \vspace{0.15em}
    \item Prosedur memakai CALL/RET dan menekankan ukuran program.

    \vspace{0.15em}
    \item Macro cocok untuk bagian pendek yang sering dipakai.
\end{itemize}
\end{frame}

% =========================================================
% CONTOH PROGRAM
% =========================================================
\begin{frame}[fragile]{Contoh 1: Prosedur Berparameter Stack}
\begin{columns}[t]
\begin{column}{0.44\textwidth}
{\color{obeTeal}\bfseries Pemanggil}\\[0.15em]
\begin{kodeasm}[numbers=none]
PUSH tinggi
PUSH lebar
CALL HitungLuas
MOV luas, AX
MOV BX, AX
\end{kodeasm}
\end{column}
\begin{column}{0.52\textwidth}
{\color{obeTeal}\bfseries Prosedur}\\[0.15em]
\begin{kodeasm}[numbers=none]
HitungLuas PROC
    PUSH BP
    MOV BP, SP
    MOV AX, [BP+4]
    MUL WORD PTR [BP+6]
    POP BP
    RET 4
HitungLuas ENDP
\end{kodeasm}
\end{column}
\end{columns}

\vspace{0.2em}
{\footnotesize\color{obeGray}Hasil $5 \times 10 = 50$ atau
\hex{0032}, disimpan di \code{luas} dan AX.\par}
\Sumber{Bab XIV, Contoh 1}
\end{frame}

% ---------------------------------------------------------
\begin{frame}[fragile]{Contoh 2: Macro dengan Parameter}
\begin{columns}[t]
\begin{column}{0.55\textwidth}
\begin{kodeasm}
Cetak_Kar MACRO Kar
    MOV CX, 3
    MOV AH, 02h
    MOV DL, Kar
Ulang:
    INT 21h
    LOOP Ulang
ENDM
.CODE
START:
    Cetak_Kar 'S'
    MOV AH, 4Ch
    INT 21h
END START
\end{kodeasm}
\end{column}
\begin{column}{0.42\textwidth}
\footnotesize
{\color{obeTeal}\bfseries Keluaran}\\[0.15em]
\begin{keluaran}
SSS
\end{keluaran}

\vspace{0.3em}
Assembler menyalin isi macro dengan setiap \code{Kar} diganti
\code{'S'}. Label \code{Ulang} aman karena macro dipakai sekali.
\end{column}
\end{columns}
\Sumber{Bab XIV, Contoh 2}
\end{frame}

% =========================================================
% PRAKTIKUM
% =========================================================
\begin{frame}{Praktikum: Stack, Prosedur, dan Macro}
\begin{block}{Judul praktikum}
\small Menelusuri stack, prosedur berparameter, dan macro\\[0.15em]
\textbf{Sub-CPMK:} 3.5 \quad\textbullet\quad
\textbf{Waktu:} 100 menit \quad\textbullet\quad
\textbf{Bentuk:} individu
\end{block}

\vspace{0.2em}
\begin{columns}[t]
\begin{column}{0.46\textwidth}
{\color{obeBlue}\bfseries Alat dan bahan}
\begin{itemize}\footnotesize
    \item Komputer dengan emu8086
    \item DOSBox dengan TASM, TLINK, dan DEBUG
    \item \file{prosedur-luas.asm} dan \file{param-bp.asm}
    \item Lembar pengamatan stack
\end{itemize}
\end{column}
\begin{column}{0.50\textwidth}
{\color{obeBlue}\bfseries Langkah kerja}
\begin{enumerate}\footnotesize
    \item Jalankan \file{prosedur-luas.asm}; catat luas.
    \item Single step; catat SP, BP, dan [BP+4]/[BP+6].
    \item Ubah parameter; pastikan luas berubah.
    \item Jalankan \file{param-bp.asm}; telusuri stack.
    \item Buat dan uji macro berparameter.
\end{enumerate}
\end{column}
\end{columns}

\vspace{0.2em}
{\footnotesize\textbf{Dikumpulkan:} berkas \file{.asm} modifikasi,
tabel SP/BP/stack, tangkapan layar, program macro, dan perbandingan
prosedur--macro.\par}
\Sumber{Bab XIV, Praktikum}
\end{frame}

% =========================================================
% LATIHAN
% =========================================================
\begin{frame}{Latihan}
\small
\begin{enumerate}
    \item Jelaskan sifat LIFO pada stack. Mengapa stack 8086 tumbuh
          ke alamat lebih rendah?

    \vspace{0.15em}
    \item Diketahui SP = \hex{FFFE}. Tentukan SP setelah PUSH AX dan
          setelah POP BX; jelaskan urutan langkahnya.

    \vspace{0.15em}
    \item Jelaskan fungsi PROC, ENDP, CALL, dan RET. Mengapa CALL
          berbeda dari JMP?

    \vspace{0.15em}
    \item Jelaskan isi [BP+0], [BP+2], [BP+4], dan [BP+6] setelah
          PUSH BP dan MOV BP, SP.

    \vspace{0.15em}
    \item Jelaskan cara mendefinisikan macro berparameter dan apa
          yang terjadi saat macro dipakai. Sebutkan empat perbedaan
          prosedur dan macro.
\end{enumerate}
\end{frame}

% =========================================================
% RANGKUMAN
% =========================================================
\begin{frame}{Rangkuman}
\small
\begin{itemize}
    \item Stack bersifat LIFO, ditunjuk SS:SP, dan tumbuh ke alamat
          rendah.

    \vspace{0.15em}
    \item PUSH mengurangi SP dua; POP menambah SP dua.

    \vspace{0.15em}
    \item Prosedur memakai PROC/ENDP; CALL menyimpan alamat kembali;
          RET memulihkannya.

    \vspace{0.15em}
    \item Stack frame dibentuk PUSH BP dan MOV BP, SP.

    \vspace{0.15em}
    \item Parameter: [BP+4] pertama, [BP+6] kedua.

    \vspace{0.15em}
    \item Macro memakai MACRO/ENDM dan parameter tekstual; isinya
          disalin saat perakitan.
\end{itemize}
\end{frame}

% ---------------------------------------------------------
\begin{frame}{Cek Pemahaman}
\footnotesize
\begin{enumerate}
    \item Setelah PUSH AX, nilai SP \dots
          \begin{enumerate}
              \item bertambah 2 \quad \item berkurang 2 \quad
              \item tidak berubah \quad \item bertambah 4
          \end{enumerate}

    \vspace{0.2em}
    \item Alamat kembali pada pemanggilan prosedur disimpan oleh
          instruksi \dots
          \begin{enumerate}
              \item JMP \quad \item CALL \quad
              \item RET \quad \item PUSH
          \end{enumerate}

    \vspace{0.2em}
    \item Setelah PUSH BP dan MOV BP, SP, parameter pertama diakses
          melalui \dots
          \begin{enumerate}
              \item [BP+0] \quad \item [BP+2] \quad
              \item [BP+4] \quad \item [BP+6]
          \end{enumerate}
\end{enumerate}
\end{frame}

% =========================================================
% PENUTUP
% =========================================================
\begin{frame}{Penutup}
\begin{block}{Sebelum pertemuan berikutnya}
\small
\begin{itemize}
    \item Selesaikan praktikum: berkas \file{.asm} modifikasi,
          tabel SP/BP/stack, program macro, dan perbandingan.
    \item Baca \textbf{Bab XV} — Array, String, dan Pencetakan
          Bilangan.
    \item Siapkan pertanyaan tentang offset array, string, dan
          konversi bilangan ke ASCII.
\end{itemize}
\end{block}

\vspace{0.4em}
\begin{center}
{\color{obeBlue}\sffamily\large\bfseries Terima kasih\par}
\vspace{0.15em}
{\color{obeGray}\sffamily\small \NamaMK\ (\KodeMK)
\textbullet\ \SemesterTeks\par}
\end{center}
\end{frame}

\end{document}
